\documentclass[11pt,a4paper]{article}

\usepackage{fontspec}
\usepackage{polyglossia}
\setmainlanguage{russian}
\setotherlanguage{english}

\IfFontExistsTF{Times New Roman}
{\setmainfont{Times New Roman}}
{\setmainfont{Libertinus Serif}}

\usepackage[a4paper,top=1.8cm,bottom=1.8cm,left=2.0cm,right=2.0cm]{geometry}
\usepackage{enumitem}
\usepackage{xcolor}
\usepackage{hyperref}
\usepackage{titlesec}

\hypersetup{colorlinks=true, linkcolor=black, urlcolor=blue}

\definecolor{hseblue}{HTML}{144F7B}
\definecolor{softblue}{HTML}{EAF4FA}
\definecolor{darktext}{HTML}{1A1A2E}

\setlength{\parindent}{0pt}
\setlength{\parskip}{0.55em}
\linespread{1.05}\selectfont

\titleformat{\section}
  {\Large\bfseries\color{hseblue}}
  {Слайд \thesection.}
  {0.5em}
  {}

\titleformat{\subsection}
  {\bfseries\color{hseblue}}
  {}
  {0pt}
  {}

\setlist[itemize]{leftmargin=1.5em, topsep=2pt, itemsep=2pt}

\newcommand{\slidetitle}[1]{\section{#1}}
\newcommand{\note}[1]{%
  \vspace{0.15em}
  \noindent\colorbox{softblue}{%
    \parbox{\dimexpr\linewidth-2\fboxsep\relax}{\textbf{Акцент:} #1}%
  }\vspace{0.25em}
}

\begin{document}

\begin{center}
{\Large\bfseries Текст выступления к защите курсовой работы\par}
\vspace{0.4em}
{\large «Исследование устойчивости моделей волатильности и ценообразования опционов в различных рыночных режимах»\par}
\vspace{0.6em}
Ковалёва Мария Сергеевна \\
НИУ ВШЭ, Факультет экономических наук, Москва, 2026
\end{center}

\vspace{0.8em}

Этот текст написан как связный сценарий выступления по финальной презентации. Его можно читать почти дословно: каждый раздел соответствует одному слайду, а таблицы и графики проговариваются словами.

\tableofcontents
\newpage

\slidetitle{Титульный слайд}

Здравствуйте, уважаемые члены комиссии. Меня зовут Ковалёва Мария Сергеевна. Тема моей курсовой работы --- «Исследование устойчивости моделей волатильности и ценообразования опционов в различных рыночных режимах».

В работе я рассматриваю, насколько точно классические модели оценивают реальные рыночные цены опционов на российском рынке. Фокус исследования --- не только на том, какая модель в среднем ближе к рынку, но и на том, как меняется ошибка модели в зависимости от состояния рынка, срока до экспирации и денежности опциона.

Практически это важно, потому что опционы используются для хеджирования, спекулятивной торговли, управления рисками и оценки рыночных ожиданий. Если модель систематически ошибается, то участник рынка может неверно оценивать стоимость риска или некорректно сравнивать разные контракты.

\note{Сразу обозначить, что работа эмпирическая: сравниваются модельные и рыночные цены на реальных данных Московской биржи.}

\slidetitle{Контекст и проблема}

На этом слайде показана логика исследовательской проблемы. На рынке опционы торгуются на бирже, и для каждого контракта у нас есть наблюдаемая рыночная цена. В моей работе в качестве рыночной цены используется расчётная цена Московской биржи.

С другой стороны, в финансовой теории и на практике для оценки «справедливой» цены опциона используются модели. В работе рассматриваются прежде всего биномиальная модель Кокса--Росса--Рубинштейна для американских опционов и модель Black-76 как европейский benchmark для опционов на фьючерсы.

Проблема в том, что любая классическая модель учитывает ограниченный набор факторов: цену базового актива, страйк, срок до экспирации, ставку, волатильность и тип опциона. Реальная рыночная цена агрегирует гораздо больше информации: ожидания участников, ликвидность, дисбаланс спроса и предложения, режим рынка, новости и другие ненаблюдаемые факторы.

Поэтому главный вопрос звучит так: насколько точны популярные модели на реальных российских данных и в каких условиях их ошибка становится больше? Иными словами, я не предполагаю, что модель должна идеально восстановить рынок. Я исследую устойчивость ошибки: где модель работает приемлемо, а где начинает систематически отклоняться.

\note{Здесь важно не обещать поиск «истинной цены»: работа сравнивает модельную цену с биржевой расчётной ценой как эмпирическим бенчмарком.}

\slidetitle{Гипотеза и цель исследования}

Гипотеза исследования состоит в том, что ошибка моделей зависит от рыночного режима и возрастает в стрессовом режиме, то есть в периоды повышенной волатильности.

Цель работы --- исследовать устойчивость моделей волатильности и ценообразования опционов в различных рыночных режимах на основе сравнения модельных и рыночных цен.

Объект исследования --- процессы ценообразования опционов на российском финансовом рынке. Предмет исследования --- устойчивость моделей волатильности и ценообразования, а также зависимость ошибки от рыночного режима, срока до экспирации, денежности опциона и способа задания волатильности.

Для достижения цели в работе решаются несколько задач. Во-первых, я рассматриваю теоретические основы моделей ценообразования и оценки волатильности. Во-вторых, формирую выборку дневных наблюдений по опционам на фьючерс MXI. В-третьих, рассчитываю модельные цены по нескольким спецификациям. И наконец, сравниваю ошибки моделей на полной выборке и в разрезах по режиму волатильности, сроку до экспирации и денежности, а также провожу формальные статистические проверки.

Ключевая идея слайда в том, что проверяется не абсолютная «правильность» теоретической цены, а устойчивость ошибки: как она устроена, где усиливается и какие факторы её объясняют.

\slidetitle{Методология исследования}

На слайде показан общий pipeline исследования.

Сначала используются данные Московской биржи. В выборку входят опционы серии MM на фьючерс MXI, то есть на фьючерс на Индекс МосБиржи. Период наблюдений --- с 3 января 2024 года по 19 мая 2026 года. Исходная база содержит 19 574 строки, 92 уникальных опционных инструмента и 5 базовых фьючерсных контрактов. После операционной фильтрации остаётся 19 565 наблюдений, то есть практически вся исходная база. Для основного сравнения трёх американских моделей используется общая сопоставимая выборка из 18 058 наблюдений, где одновременно доступны все оценки волатильности.

Дальше для каждого базового фьючерса рассчитывается волатильность тремя способами. Первый способ --- историческая волатильность за 21 торговый день, то есть более короткое и быстро реагирующее окно. Второй способ --- историческая волатильность за 63 торговых дня, то есть более устойчивое квартальное окно. Третий способ --- прогнозная условная волатильность GARCH(1,1), которая оценивается на расширяющемся окне без look-ahead bias.

Затем эти оценки волатильности подставляются в модели ценообразования. Основная модель --- CRR American, то есть биномиальная модель для американского опциона. Отдельно рассчитывается Black-76 как европейский benchmark без возможности досрочного исполнения.

После этого для каждого наблюдения считается ошибка как «модельная цена минус рыночная цена». Если ошибка отрицательная, модель недооценивает рынок; если положительная --- переоценивает.

Качество оценивается по стандартным метрикам MAE, RMSE и средней ошибке. Кроме того, в работе вводится авторская метрика TVNAE. Она нормирует абсолютную ошибку на временну́ю стоимость опциона, то есть на ту часть цены, которая действительно зависит от волатильности, срока и ставки. Это особенно важно для глубоких опционов в деньгах, где большая часть цены --- внутренняя стоимость.

\note{При описании схемы удобно вести комиссию сверху вниз: данные → волатильность → модель → ошибка → разрезы анализа.}

\slidetitle{Модели ценообразования и входные параметры}

На этом слайде я подробнее показываю две модели ценообразования.

Первая модель --- CRR American, или биномиальная модель Кокса--Росса--Рубинштейна. Она строит дерево возможных цен базового актива и на каждом узле позволяет сравнить стоимость продолжения с немедленным исполнением. Именно поэтому она подходит для американских опционов: держатель такого опциона может исполнить контракт досрочно в любой момент до экспирации.

Вторая модель --- Black-76. Это аналитическая европейская модель для опционов на фьючерсы. Она удобна и широко используется как benchmark, но не учитывает право досрочного исполнения. Поэтому в моей работе Black-76 не является полноценной конкурирующей моделью для американских опционов, а используется как эталон без раннего исполнения, чтобы оценить, насколько важен сам выбор класса модели.

Входные параметры у моделей похожи: цена фьючерса, страйк, срок до экспирации, безрисковая ставка, волатильность и тип опциона. В базовой спецификации ставка фиксируется на уровне 16 процентов годовых, что соответствует ключевой ставке ЦБ в рассматриваемой постановке.

Главный варьируемый параметр --- волатильность. В работе сравниваются HV\_21d, HV\_63d и GARCH. Идея такая: разные оценки волатильности дают разные модельные цены, а значит и разные ошибки. Поэтому по сути сравнивается не только модель ценообразования, но и связка «модель цены плюс модель волатильности».

\slidetitle{Общее сравнение моделей}

Здесь представлены основные результаты по полной сопоставимой выборке для трёх американских спецификаций.

В таблице видно, что по MAE лидирует CRR с GARCH-волатильностью: средняя абсолютная ошибка равна 107.0 пункта. Очень близкий результат у CRR с 63-дневной исторической волатильностью: 107.1 пункта. Разница между ними практически нулевая. CRR с 21-дневной исторической волатильностью хуже: MAE равен 119.3 пункта.

Если смотреть на RMSE, то лучшей становится 63-дневная историческая волатильность: RMSE равен 140.6 пункта против 145.8 у GARCH. Это означает, что HV\_63d немного устойчивее к крупным промахам. Но по типичной абсолютной ошибке GARCH всё же немного лучше.

Важнейший результат таблицы --- средняя ошибка отрицательная у всех моделей. У GARCH она равна минус 78.1 пункта, у HV\_63d --- минус 89.7, у HV\_21d --- минус 95.3. Это значит, что все три американские спецификации систематически недооценивают рыночные расчётные цены.

График в нижней части слайда визуально подтверждает этот вывод. На нём сравниваются ошибки моделей: GARCH и HV\_63d расположены очень близко по MAE, 21-дневная историческая волатильность заметно отстаёт, а отрицательное смещение является общей проблемой всех спецификаций.

Интерпретация такая: простая постановка с единственной оценкой волатильности на контракт не устраняет систематическое смещение. Скорее всего, рынку нужна более гибкая структура, например поверхность подразумеваемой волатильности, а не одна плоская sigma.

\slidetitle{Почему стандартных метрик недостаточно: TVNAE}

Этот слайд объясняет, зачем в работе вводится дополнительная метрика TVNAE.

Стандартные метрики, такие как MAE и RMSE, считают ошибку по полной цене опциона. Но в моей выборке много длинных опционов глубоко в деньгах. У таких контрактов большая часть цены --- это внутренняя стоимость, то есть механическая разница между ценой фьючерса и страйком. Эту часть все модели воспроизводят примерно одинаково, потому что она почти не зависит от качества прогноза волатильности.

Поэтому абсолютная ошибка может вводить в заблуждение. Например, ошибка в 100 пунктов выглядит большой, но если временна́я стоимость опциона равна 500 пунктам, это 20 процентов от опциональной компоненты. А если временна́я стоимость равна 80 пунктам, то те же 100 пунктов означают, что модель ошиблась больше чем на всю временную стоимость.

Именно поэтому я использую TVNAE: абсолютная ошибка делится на временну́ю стоимость опциона. Временная стоимость определяется как рыночная цена минус внутренняя стоимость.

В таблице видно, что по TVNAE GARCH уже имеет более заметное преимущество. Среднее TVNAE у GARCH равно 0.66, у HV\_63d --- 0.70, у HV\_21d --- 0.76. Медианное TVNAE также лучше у GARCH: 0.62 против 0.68 и 0.78.

График-heatmap на слайде показывает TVNAE по моделям и срокам до экспирации. Самый сильный результат --- на коротких опционах. Для DTE 0--7 дней TVNAE у GARCH равно 0.19, а у HV\_63d --- 0.59. То есть по временной стоимости GARCH примерно в три раза точнее на самых коротких контрактах.

Главный вывод этого слайда: стандартные метрики почти уравнивают GARCH и HV\_63d, потому что выборка смещена в сторону длинных и deep ITM опционов. TVNAE показывает более содержательную картину: GARCH лучше описывает именно опциональную, временную компоненту цены.

\slidetitle{Американская модель vs Black-76}

На этом слайде сравнивается американская модель CRR и европейская модель Black-76 при одинаковой 63-дневной исторической волатильности.

Результат очень показательный. У американской модели CRR + HV\_63d MAE равен 107.1 пункта, RMSE --- 140.6, средняя ошибка --- минус 89.7. У Black-76 + HV\_63d MAE уже 156.2 пункта, RMSE --- 211.6, средняя ошибка --- минус 149.0. То есть европейская модель хуже по MAE примерно на 46 процентов и имеет намного более сильное отрицательное смещение.

Причина не в том, что 63-дневная волатильность вдруг стала плохой. Причина в классе модели: Black-76 не учитывает досрочное исполнение. А рассматриваемые опционы являются американскими, и держатель имеет право исполнить их до экспирации.

На слайде отдельно выделен факт: в 24.5 процента наблюдений цена Black-76 оказалась ниже внутренней стоимости опциона. Это уже не просто неточность, а нарушение естественной нижней границы стоимости опциона. Американский опцион не должен стоить меньше того, что можно получить при немедленном исполнении.

Нижний график-heatmap показывает премию раннего исполнения по режимам волатильности и срокам до экспирации. Средняя премия раннего исполнения по всей выборке равна 59.28 пункта, медиана --- 9.87. На коротких сроках премия почти нулевая, но для длинных контрактов резко возрастает. Особенно сильно это видно в режиме высокой волатильности при DTE 91+: средняя премия составляет примерно 106 пунктов.

Ключевой вывод: для этих контрактов выбор класса модели важнее, чем выбор оценки волатильности. Если опцион американский, европейская формула систематически занижает цену, особенно на длинных контрактах и в стрессовых условиях.

\slidetitle{Ошибки по рыночным режимам}

На этом слайде показано, как средняя абсолютная ошибка меняется по режимам волатильности.

В режиме низкой волатильности лучшей оказывается GARCH-спецификация: MAE равен 77.5 пункта. HV\_63d даёт 82.1 пункта, а HV\_21d заметно хуже --- 104.1 пункта. В среднем режиме GARCH также лидирует: 118.2 пункта против 128.3 у HV\_63d и 141.8 у HV\_21d.

Но в режиме высокой волатильности картина меняется. GARCH ухудшается до 141.0 пункта, а исторические оценки становятся лучше: 122.8 у HV\_63d и 117.8 у HV\_21d. Это важный и не совсем очевидный результат.

Возможная интерпретация такая: GARCH хорошо работает в спокойных и умеренных условиях, потому что улавливает текущий уровень условной волатильности. Но в стрессовом режиме модель может быть инерционной: она либо медленно перестраивается при резком изменении режима, либо переоценивает персистентность шоков. Поэтому её преимущество исчезает.

График rolling MAE в нижней части слайда показывает, что ошибка моделей во времени не стабильна: есть периоды, когда разрыв между спецификациями расширяется, и периоды, когда они сближаются. Это поддерживает основную идею исследования: качество модели зависит от рыночного контекста.

При этом важно уточнение из регрессионного анализа: режим волатильности влияет на ошибку описательно, но не является единственным структурным фактором. Часть эффекта высокой волатильности объясняется тем, что в этом режиме в выборке много длинных контрактов и определённых зон денежности.

\slidetitle{Ошибки по сроку до экспирации и денежности}

Здесь я показываю один из самых важных разрезов --- ошибки по сроку до экспирации.

На коротких опционах с DTE 0--7 дней все ошибки малы, но GARCH всё равно лучше: MAE 1.88 пункта против 3.53 у HV\_63d и 2.78 у HV\_21d. На сроке 8--30 дней преимущество GARCH становится ещё более заметным: 4.87 пункта против примерно 10 пунктов у HV\_63d и 8.7--8.8 у HV\_21d. На сроке 31--90 дней GARCH также выигрывает: 24.6 пункта против 35--36 пунктов у исторических оценок.

А вот на длинных контрактах, то есть DTE 91+, преимущество GARCH исчезает. Здесь HV\_63d немного лучше: 121.0 пункта против 122.5 у GARCH. HV\_21d заметно хуже --- 135.3 пункта.

Экономический смысл такой: GARCH даёт однодневный прогноз условной волатильности, поэтому он хорошо отражает текущую турбулентность на коротком горизонте. Но для длинных опционов цена зависит от ожиданий волатильности на месяцы вперёд. Одного краткосрочного прогноза GARCH уже недостаточно.

Слайд также отсылает к денежности. По денежности GARCH в отчёте лидирует почти во всех зонах: deep OTM, OTM, ATM и ITM. Исключение --- deep ITM, где 63-дневная историческая волатильность немного лучше. Крупные абсолютные ошибки deep ITM частично связаны с тем, что цена таких опционов содержит большую внутреннюю стоимость. Поэтому, опять же, важно смотреть не только MAE, но и TVNAE.

График на слайде визуализирует разложение ошибок: он помогает увидеть, что главная масса ошибки концентрируется в длинных контрактах, а не в коротких.

\slidetitle{Формальные проверки и устойчивость}

На этом слайде собраны формальные проверки результатов.

Первая проверка --- тест средней ошибки. Для каждой американской модели проверяется гипотеза, что средняя ошибка равна нулю. Стандартные ошибки кластеризуются по торговой дате, чтобы учесть зависимость наблюдений внутри одного дня. Результат однозначный: у всех трёх моделей t-stat находится примерно от минус 20 до минус 25, p-value меньше 0.001. Значит, систематическая недооценка статистически значима.

Вторая проверка --- тесты Диболда--Мариано для попарного сравнения функций потерь на дневном уровне. По ним GARCH и HV\_63d статистически неразличимы на полной выборке: p-value по MAE около 0.770, по MSE около 0.370. При этом HV\_21d статистически значимо хуже HV\_63d: p-value около 0.0003 по MAE.

Третья часть --- описательная OLS-регрессия абсолютной ошибки на характеристики наблюдения. В финальной спецификации R-квадрат равен примерно 0.27. Это значит, что срок до экспирации, денежность, тип опциона и режим объясняют около 27 процентов вариации ошибки. Коэффициент при DTE положительный: каждый дополнительный день до экспирации увеличивает абсолютную ошибку примерно на 0.16 пункта. Коэффициент при абсолютной логарифмической денежности отрицательный: около денег ошибки выше, а в хвостах ниже после контроля прочих факторов. Путы оцениваются точнее коллов: коэффициент около минус 28 пунктов.

Важное уточнение: после добавления DTE и денежности dummy высокой волатильности становится отрицательным. Это означает, что сам по себе high-vol regime не является главным источником роста ошибки. В простых разрезах он выглядит важным, потому что в высокой волатильности меняется состав наблюдений, в частности больше длинных контрактов.

Справа показаны robustness checks. На полной выборке GARCH и HV\_63d очень близки. Без deep OTM и deep ITM лидирует GARCH. На DTE до 91 дня GARCH явно лучше. На DTE больше 91 дня HV\_63d и GARCH практически равны. Без high-vol режима снова лидирует GARCH. При изменении ставки на 12, 16 и 20 процентов качественные выводы сохраняются.

\note{Этот слайд закрывает вопрос статистической значимости: результаты не только описательные, но и проверены формально.}

\slidetitle{Итоговые выводы}

На этом слайде я формулирую основные выводы работы.

Первый вывод: гипотеза подтверждается в уточнённой форме. Ошибка действительно зависит от рыночного режима, но многофакторный анализ показывает, что значительная часть этого эффекта объясняется сроком до экспирации и денежностью. То есть stress regime важен, но не сам по себе: важно, какие именно контракты преобладают в этом режиме.

Второй вывод: для рассматриваемых контрактов нужна американская модель. Black-76 даёт MAE на 46 процентов хуже и в 24.5 процента наблюдений оценивает опцион ниже внутренней стоимости. Это показывает, что право досрочного исполнения является содержательным фактором, а не технической деталью.

Третий вывод: 63-дневная историческая волатильность --- устойчивый и простой baseline. На полной выборке она статистически неотличима от GARCH, а по RMSE даже немного лучше, то есть лучше контролирует крупные промахи.

Четвёртый вывод: GARCH полезен прежде всего там, где важна временная стоимость и короткий горизонт. По TVNAE он лучше: 0.66 против 0.70 у HV\_63d. На DTE 0--7 дней преимущество особенно сильное: TVNAE 0.19 против 0.59, то есть примерно в три раза меньше ошибка по временной стоимости.

Пятый вывод: все модели систематически недооценивают рынок. Это значит, что простая постановка с одной волатильностью не способна полностью описать реальные цены. Для дальнейшего улучшения нужны более гибкие подходы: поверхность подразумеваемой волатильности, учёт ликвидности, модели со стохастической волатильностью и скачками.

\slidetitle{Ограничения и направления развития}

Здесь я отдельно фиксирую ограничения исследования.

Первое ограничение --- в качестве рыночной цены используется расчётная цена биржи, а не котировка или цена последней сделки. Для неликвидных инструментов эти значения могут расходиться, поэтому расчётная цена является удобным и воспроизводимым бенчмарком, но не абсолютной истинной стоимостью.

Второе ограничение --- структура выборки. Она смещена в сторону длинных опционов в деньгах. Это важно для интерпретации результатов: выводы хорошо описывают именно такую структуру рынка, но их нужно осторожно переносить на рынки с другой структурой контрактов.

Третье ограничение --- единая оценка волатильности без поверхности. Каждая модель получает одну sigma на контракт, но реальный рынок обычно имеет улыбку и временную структуру подразумеваемой волатильности. Поэтому систематическое смещение в этой постановке полностью устранить нельзя.

Четвёртое ограничение --- фиксированная ставка. В базовой спецификации используется 16 процентов годовых, но реальная срочная структура ставок может быть неоднородной. В проверках чувствительности изменение ставки не меняет ранжирование моделей, однако более точная ставка могла бы снизить часть смещения.

Направления развития напрямую следуют из этих ограничений. Первое --- переход к поверхности подразумеваемой волатильности. Второе --- учёт ликвидностных факторов, потому что расчётные цены могут включать премию за неликвидность. Третье --- применение моделей со стохастической волатильностью, например Heston или SABR. Четвёртое --- модели со скачками, поскольку доходности финансовых активов часто имеют тяжёлые хвосты. Пятое --- использование срочной структуры ставок вместо одного фиксированного значения.

В целом дальнейшая логика исследования --- перейти от простых benchmark-моделей к более гибким моделям, которые лучше отражают реальную структуру рынка.

\slidetitle{Спасибо за внимание}

Спасибо за внимание. Я готова ответить на вопросы.

Если понадобится коротко повторить главный результат, его можно сформулировать так: на российских данных по опционам на фьючерс MXI американская постановка принципиально важна, GARCH даёт преимущество на коротких и среднесрочных контрактах и по временной стоимости, но на полной выборке статистически не отличается от устойчивой 63-дневной исторической волатильности; при этом все простые модели систематически недооценивают рынок, что указывает на необходимость более гибкой волатильностной структуры.

\end{document}
